% ==============================================================================
%  CHAPTER 5 — DISCUSSION
% ==============================================================================
\section{Discussion}
\label{sec:discussion}

\subsection{What the Accuracy Results Show}
\label{sec:disc_perf}

The clearest result is that the full system errs on the side of privacy. In the video-clip
study it left only one bystander visible out of 33, and when it did keep a face visible
the person was almost always a partner. The cost is on the other side: it blurred about
four in ten partners. The still-frame study showed the same pattern on a smaller sample,
and because the two studies agree, I read this as a property of the system rather than an
accident of one sample.

I expected the landmark cues to change both kinds of error: the head-pose veto should
remove bystanders who happen to be close and central, and the speaking cue should rescue
partners who are slightly off-centre. The first expectation held, since almost no
bystanders were exposed. The second did not, but not because the landmark cues failed.
When I examined the missed partners, most of them never reached the speaking cue at all:
their faces were smaller than the proximity threshold, and the speaking cue only applies to
faces that pass that threshold. The typical case was a friend sitting on the far side of
a card table, whose face covered less than 1\% of the frame. The 1.2\% threshold separates
close conversation partners from distant bystanders well, but a group around a table puts
many partners in between.

The second cause was timing. Ten of the missed partners were confirmed by the system
later in the same video, after their clip had ended. The live system needs the cues to
hold for two consecutive frames before confirming anyone, and a friend who is small or
partly turned may take a while to meet them. This is what motivated the offline second
pass, and it explains why simply applying the latch backwards in time recovered so many
partners.

Comparing the two phases is harder than the numbers in Table~\ref{tab:vs_baseline}
suggest. Phase~2a was judged by me alone, on a different sample and with a partner
definition of its own (a face-to-face conversation at any point in the video), and the
confidence intervals of the two phases overlap. What can be said is that adding the landmark cues moved the system
towards privacy, and that the remaining partner losses come from geometry and timing.

\subsection{Comparison with BystandAR}
\label{sec:disc_bystandar}

BystandAR protected 98.14\% of bystanders while keeping 96.27\% of intended subjects
visible \cite{bystandar}. My live system matches its bystander protection closely, but
keeps far fewer partners. Two differences explain most of this. BystandAR knows where the
wearer is looking from an eye tracker, which tells it directly who the wearer is engaging
with; my system has to infer this from where faces appear in the frame, and a friend
sitting off to one side looks, geometrically, much like a stranger. BystandAR was also
evaluated in controlled sessions with planned interactions, while my clips are
unscripted videos of card games and shopping. Its partners stand where a conversation
partner normally stands; many of mine sit across a table.

The point of my approach is that it needs neither the eye tracker nor the microphone
array. Glasses such as the Ray-Ban Meta have a camera but no eye tracking, so BystandAR's
method cannot run on them, while mine can.

\subsection{Is the System Practical?}
\label{sec:disc_fps}

Whether the measured speed is practical depends on how the system is used.

For \emph{recorded video}, it is practical now. A five-minute recording is processed on
the laptop in about 15 minutes, and the offline second pass, which gives the best
accuracy, only makes sense for recordings anyway. This matches how blanket anonymization
is used today: EgoBlur, for example, is applied to recordings from Meta's Aria glasses
before the data is shared \cite{egoblur2023}.

For \emph{live use on glasses}, the current implementation is too slow. The projected
1.5 to 1.6 frames per second on a Raspberry Pi~5 (Table~\ref{tab:projection}) would not
keep up with a live camera. The breakdown in Table~\ref{tab:fps} shows where the problem
lies: detection takes almost two thirds of the time, while the classification method that
is the contribution of this thesis takes 0.02~ms. The bottleneck is therefore a standard
detector run in its slowest form, and Table~\ref{tab:det_hw} shows how much is available
without changing the method. Converting the same detector to NCNN makes it more than five
times faster on the Pi~5, which roughly triples the projected frame rate. Moving detection
to a neural processing unit, which every current camera-glasses chip includes, would
remove most of what remains: on a phone's NPU the detector takes 11.3~ms.

A further gain that does not depend on hardware is to run the detector only on some
frames. People in a conversation move little between consecutive frames, and the tracker
already predicts where each face will be, so the detector could run every second or third
frame with the tracker filling the gaps. I did not test this, so I cannot report how much
accuracy it would cost.

\paragraph{Comparison with commercial camera glasses.}
Table~\ref{tab:glasses} compares my evaluation hardware with current camera glasses.

\begin{table}[htbp]
\begin{singlespace}
\centering
\caption{Evaluation hardware compared with current camera glasses, from manufacturer and
press sources. Where a figure has not been released, the table says so.}
\label{tab:glasses}
\renewcommand{\arraystretch}{1.3}
\small
\setlength{\tabcolsep}{4pt}
\begin{tabular}{@{}>{\raggedright\arraybackslash}p{0.20\linewidth}
                 >{\raggedright\arraybackslash}p{0.21\linewidth}
                 >{\raggedright\arraybackslash}p{0.16\linewidth}
                 >{\raggedright\arraybackslash}p{0.15\linewidth}
                 >{\raggedright\arraybackslash}p{0.20\linewidth}@{}}
\toprule
\textbf{Device} & \textbf{Processor} & \textbf{AI accelerator} & \textbf{Battery} &
\textbf{Where vision AI runs} \\
\midrule
Laptop (my evaluation) & Intel Core i5-13420H & none used & mains or battery & on the device (CPU) \\
Raspberry Pi 5 (projection target) & 4$\times$ Cortex-A76, 2.4~GHz & none & external & on the device (CPU) \\
Ray-Ban Meta \cite{qcom_ar1} & Snapdragon AR1 Gen~1 & Hexagon NPU & about 4~h & mostly phone and cloud \\
Meta Ray-Ban Display \cite{gsmarena_display} & Snapdragon AR1 Gen~1, 2~GB RAM & Hexagon NPU & about 6~h & mostly phone and cloud \\
Snap Spectacles, 5th gen.\ \cite{spectacles5} & 2$\times$ Snapdragon & not published & about 45~min & on the device \\
Meta Aria Gen~2 \cite{aria2} & custom co-processor & custom & 6-8~h & on the device \\
\bottomrule
\end{tabular}
\end{singlespace}
\end{table}

The comparison supports three points. First, the Pi~5 is a reasonable stand-in for the CPU
of these glasses: its four Cortex-A76 cores at 2.4~GHz are of the same class as the AR1's
four cores at 1.9~GHz, which is why the Ray-Ban Meta projection in
Table~\ref{tab:projection} sits slightly below the Pi~5. Second, every glasses platform
has a neural processing unit and my evaluation used none, so the CPU-only figures are a
worst case. Third, battery capacity is a tighter limit than speed. Spectacles run for
about 45 minutes, and the Ray-Ban Meta sends most of its vision processing to the phone or the cloud. Running my pipeline continuously would compete for the
same battery, which makes running the detector less often even more important. Keeping the
processing on the glasses has a privacy benefit of its own: bystanders' faces never leave
the device, which Meta itself identifies as a central concern for camera glasses
\cite{meta_bystander}.

\subsection{Failure Analysis}
\label{sec:disc_failures}

\paragraph{Phase 2a.}
All five bystanders that Phase~2a left visible came from interaction videos: people who
stood close to the wearer and near the centre of the view for a long time without ever
being addressed. Proximity and centrality cannot
separate them from partners, which is the situation the head-pose veto was added for. The
four partners Phase~2a missed had short interactions after spending a long time at the
edge of the frame, so the two-frame confirmation was not reached before the interaction
ended.

\paragraph{Phase 2b.}
All three partners missed in the still-frame study had faces covering 0.36\% to 0.72\% of
the frame, below the 1.2\% threshold. Two of these were the same person, covered by two
overlapping detections that SORT tracked separately, so one miss was counted twice. The
third stood at a shop counter near the edge of the frame. In the video-clip study the
pattern repeated on a larger scale: 25 of the 30 missed partners were below the area
threshold, with a median face size of 0.91\% of the frame against 1.72\% for the partners
the system kept.

The landmark cues were also available less often than I had assumed. MediaPipe found a
face in only 50.8\% of the 224{,}519 face crops the pipeline logged. Most failures were on
small, blurred or partly hidden faces, which are exactly the distant partners the
speaking cue was meant to rescue.

\paragraph{Detector errors.}
Two of the 104 sampled tracks were not people. One was a face printed on a magazine
cover, which the system blurred. The other was a product package on a shop shelf, which
the system confirmed as a partner and left unblurred. Neither exposes a real person, but
the second shows that the cues are applied to anything the detector reports. That
detection had a median confidence of 0.53, just above the threshold of 0.45, and MediaPipe
found a face in only 1 of its 73 frames, so requiring successful landmark extraction before
confirming a partner would have rejected it.

\subsection{Limitations}
\label{sec:disc_limitations}

\paragraph{The dataset.}
All 21 videos came from public online sources, so I could not control the camera, its
mounting or the frame rate, and I do not know who recorded them or why. The two categories
also differ in setting: every interaction video is indoors, mostly at a table, and every
bystander video is outdoors. A conversation in the street, or a crowded room in which the
wearer is talking to no one, is not represented, and the system may behave differently
there.

\paragraph{The labels.}
I labelled Set~A alone, and I was one of the annotators in Sets B and C while knowing how
the system worked. In Set~C the other four annotators agreed with my labels on 92\% to 98\%
of clips, and removing my labels does not change the majority label of any clip, but an
evaluation without my involvement would still be cleaner. Set~B and Set~C also used different
partner definitions, so their results are not strictly comparable.

\paragraph{Tuning on the evaluation data.}
I chose the classification thresholds by looking at my own videos, and I designed the
second-pass rules after seeing the video-clip errors, then measured them on the same
clips. The sensitivity check on the group rule's time limit makes it unlikely that the
exact value was overfitted, but only a fresh, separately labelled set of clips can show
how well the thresholds and rules generalise.

\paragraph{Speech without sound.}
The speaking cue only sees whether the mouth is open. Laughing, eating or chewing look the
same as speaking, and a partner who listens silently gets no help from this cue at all.

\paragraph{Speed and energy.}
The Raspberry Pi~5 and glasses figures are projections built from published benchmarks
and stated assumptions, not measurements on those devices, and the battery-power runs
were affected by the processor heating up. % TODO(Akash): replace this sentence once the
% power measurement has been run.
I have not measured how much energy the pipeline uses, which matters at least as much as
speed for a device with a 45-minute battery.
